\documentclass{subfiles}
\begin{document}\label{sec:RB-insert}
    Inserting a node in a RB tree works similarly to that in a BST, 
        with a couple of differences. First, we color the node just added red,
        then, to ensure all RB properties hold a auxiliary procedure is executed.
        \autoref{proc:RB-insert} is the pseudo code used to implement insertion in a red-black tree,
        while \autoref{proc:RB-insert-fixup} the pseudo code for the auxiliary function is given.

    Before analyzing how \autoref{proc:RB-insert-fixup} works, let us see which of the RB properties
        may be violated. Properties 1, 3, 5 still hold\footnotemark.
        On the other hand, properties 2 and 4 are not garanteed to hold; 
        the former failing when the new node becomes the root, the latter if the parent is also red.

    \subfile{../TikZ/Procedure 2 - RB insert}

    The analyzing of \autoref{proc:RB-insert-fixup} can be broke down in six cases, 
        pairwise symmetric. We consider for this reason case 1-3.
        Let \(z\) be the inserted node, let \(z.p\) denote \(z\)'s parent and let \(y\) denote \(z\)'s uncle.
        \begin{description}
            \item [Case 1: both \(\mathbf{z.p}\) and \(\mathbf{y}\) are red.] 
                Since both \(z.p\) and \(y\) are red, 
                and we can assume \(z.p.p\) to exist and being black,
                we simply recolor the nodes, locally restoring property 4.
                Then, we call the auxiliary procedure on \(z.p.p\).

            \item [Case 2: \(\mathbf{y}\) is black and \(\mathbf{z}\) is a right child.]
                We apply a left rotation on \(z.p\) falling this way in case 3.

            \item [Case 3: \(\mathbf{y}\) is black and \(\mathbf{z}\) is a left child.] 
                Since \(z.p.p\) was black prior to the insertion,
                we can simply apply some recolorig a right rotation on \(z.p.p\).
        \end{description}
        Once the loop has finished the only property that may still fail is property 2,
        which is fixed by recolorig the root black.

    \footnotetext{
        We assume that the right and left children of the inserted node are set to \textbf{nil},
        which is assumed to have a black color.
    }

    \subfile{../TikZ/Procedure 3 - RB insert fixup}
    \clearpage
\end{document}
